\documentclass[aspectratio=169,11pt]{beamer}

\usetheme{Madrid}
\usecolortheme{default}
\usefonttheme{professionalfonts}
\setbeamertemplate{navigation symbols}{}
\setbeamertemplate{footline}[frame number]

\usepackage{amsmath, amssymb, booktabs}

\title{Housing, Rents, Moving Costs}
\subtitle{And Labor-Market Adjustment}
\author{Special Topic 4: Urban and Labor -- Week 2}
\date{}

\begin{document}

\begin{frame}
  \titlepage
\end{frame}

\begin{frame}{Central question}
\begin{itemize}
    \item When do housing markets and moving costs prevent workers from reaching better labor-market opportunities?
    \item Housing is a labor-market constraint because it governs who can live near, commute to, and remain attached to jobs.
    \item The key object is real access, not nominal wages alone.
    \item The research task is to identify which friction is binding.
\end{itemize}
\end{frame}

\begin{frame}{Labor-space-housing framework}
\[
U_{ijr} = w_j - R_r - \tau(d_{jr}) + A_r - M_{ir} + \varepsilon_{ijr}
\]
\small
\begin{tabular}{p{0.22\linewidth} p{0.62\linewidth}}
\toprule
Term & Interpretation \\
\midrule
$w_j$ & wage at workplace $j$ \\
$R_r$ & rent or housing cost at residence $r$ \\
$\tau(d_{jr})$ & generalized commute cost between workplace and residence \\
$A_r$ & local amenity or disamenity value \\
$M_{ir}$ & worker-specific moving or attachment cost \\
\bottomrule
\end{tabular}
\end{frame}

\begin{frame}{How housing enters labor outcomes}
\begin{itemize}
    \item Rents convert nominal wages into real access.
    \item Housing supply elasticity determines whether shocks move prices or population.
    \item Moving costs make migration responses unequal and sometimes slow.
    \item Wealth, credit, tenure, and liquidity constraints shape who can enter high-opportunity places.
    \item Commuting can substitute for migration, but only when transport access is feasible.
    \item Local shocks are absorbed through wages, rents, employment, population, firms, and flows.
\end{itemize}
\end{frame}

\begin{frame}{Nominal wage gain versus real access gain}
\[
\Delta RA_i = \Delta w_i - \Delta R_i - \Delta \tau_i + \Delta A_i - \Delta M_i
\]
\begin{itemize}
    \item A higher wage is not automatically a higher-opportunity place.
    \item Rent capitalization can transfer part of the gain to landlords or incumbent homeowners.
    \item Commute costs can preserve access without migration, or erase the practical value of a wage gain.
    \item Worker heterogeneity matters: skill, wealth, tenure, family constraints, and credit access change the net gain.
\end{itemize}
\end{frame}

\begin{frame}{Supply constraints versus moving costs}
\small
\begin{tabular}{p{0.25\linewidth} p{0.32\linewidth} p{0.32\linewidth}}
\toprule
Friction & Signature & Counterfactual \\
\midrule
Housing supply constraint & high rents, high prices, limited population growth & allow more housing in productive places \\
Moving-cost friction & workers do not relocate despite wage gaps & reduce relocation, search, liquidity, or tenure barriers \\
Credit or wealth constraint & access differs by liquidity, ownership, or borrowing capacity & relax financing or upfront-cost barriers \\
Regulatory constraint & price adjustment concentrated where building is hard & change land-use or permitting limits \\
\bottomrule
\end{tabular}
\end{frame}

\begin{frame}{Migration and commuting}
\begin{itemize}
    \item Migration changes residence and the whole job-residence bundle.
    \item Commuting links workers to jobs without residential relocation.
    \item A local shock can increase in-commuting even when local population barely changes.
    \item Commuting openness changes local employment elasticities and incidence.
    \item Inequality appears when some workers can move, others can commute, and others can do neither.
\end{itemize}
\end{frame}

\begin{frame}{Incidence: who captures local opportunity?}
\small
\begin{tabular}{p{0.27\linewidth} p{0.55\linewidth}}
\toprule
Claim & What to check \\
\midrule
Workers gain & real wages rise after rents and commuting costs \\
Landlords gain & rents rise faster than worker real access \\
Homeowners gain & house prices capitalize the shock \\
Firms gain & labor supply expands with limited wage pressure \\
Commuters gain & workplace employment rises without resident gains \\
Would-be migrants lose & productive places price out entrants \\
\bottomrule
\end{tabular}
\end{frame}

\begin{frame}{Equilibrium adjustment}
\[
\Delta Y_\ell =
\left(
\Delta w_\ell,\,
\Delta R_\ell,\,
\Delta N_\ell,\,
\Delta E_\ell,\,
\Delta C_\ell,\,
\Delta P_\ell
\right)
\]
\begin{itemize}
    \item $w$: wages, $R$: rents, $N$: residents, $E$: employment.
    \item $C$: commuting flows, $P$: house prices or local prices.
    \item The pattern across margins reveals the likely binding friction.
    \item One coefficient rarely tells the whole incidence story.
\end{itemize}
\end{frame}

\begin{frame}{Paper architecture}
\small
\begin{tabular}{p{0.28\linewidth} p{0.55\linewidth}}
\toprule
Bucket & Research role \\
\midrule
Foundational framework & wages, rents, amenities, and location choice \\
Housing constraints and misallocation & productive places blocked by housing costs or supply \\
Moving costs and migration adjustment & workers fail to move after opportunity or decline \\
Commuting and within-metro access & workers reach jobs without changing residence \\
Local shocks and incidence & surplus divided among workers, firms, owners, and landlords \\
\bottomrule
\end{tabular}
\end{frame}

\begin{frame}{Data sources researchers use}
\small
\begin{tabular}{p{0.28\linewidth} p{0.56\linewidth}}
\toprule
Source & Typical use \\
\midrule
ACS / ACS PUMS & rents, tenure, migration, commutes, worker traits \\
Decennial Census / long form & long-run population, housing, and sorting \\
LEHD / LODES / QWI & workplace-residence links and flows \\
IRS migration & county and state mobility responses \\
Zillow / FHFA / CoreLogic & prices, rents, transactions, capitalization \\
HUD / vouchers / HMDA & assistance, mobility supports, credit access \\
CZ / MSA / tract / ZIP delineations & labor-market and housing-market geography \\
Supply / regulation measures & elasticity, zoning, permits, land constraints \\
\bottomrule
\end{tabular}
\end{frame}

\begin{frame}{Research architecture / design map}
\small
\begin{tabular}{p{0.30\linewidth} p{0.50\linewidth}}
\toprule
Design family & Best use \\
\midrule
Spatial equilibrium / counterfactual & aggregate misallocation and welfare incidence \\
Bartik-style / shift-share shock & wage, rent, employment, and population incidence \\
Housing-supply heterogeneity & price versus quantity adjustment \\
Migration event study & timing and heterogeneity of relocation \\
Commuting or transport shock & access changes without migration \\
Administrative linkage & worker, firm, residence, and workplace heterogeneity \\
\bottomrule
\end{tabular}
\end{frame}

\begin{frame}{What makes a contribution}
\begin{itemize}
    \item Identify the binding friction rather than naming housing generically.
    \item Separate nominal wage effects from real access effects.
    \item Use a geography that matches the labor-market margin.
    \item Measure both residence and workplace outcomes when possible.
    \item State incidence: workers, landlords, homeowners, firms, commuters, or excluded entrants.
\end{itemize}
\end{frame}

\begin{frame}{Discussion prompts}
\begin{itemize}
    \item A city has higher nominal wages. What data are needed before calling it a better opportunity?
    \item A demand shock raises employment and rents. Who gained?
    \item How would you distinguish a housing supply constraint from a moving-cost friction?
    \item When does transit policy create labor access without migration?
\end{itemize}
\end{frame}

\begin{frame}{Bridge to Week 3}
\begin{itemize}
    \item Week 2 asks when housing blocks or redirects access to productive places.
    \item Week 3 moves inside the metro area: spatial mismatch, segregation, neighborhoods, networks, and job access.
    \item The same discipline carries forward: define reachable jobs, housing costs, feasible commutes, and incidence.
\end{itemize}
\end{frame}

\begin{frame}{Takeaways}
\begin{itemize}
    \item Housing is part of labor-market adjustment, not background context.
    \item Real access depends on wages, rents, commuting, amenities, and moving costs.
    \item Migration and commuting are distinct adjustment margins.
    \item The best papers identify which friction is binding and who captures the surplus.
\end{itemize}
\end{frame}

\end{document}
